\scanpage{11}
\mspage{4}

Le groupe $\mathbb Z/6\mathbb Z$ opère transitivement sur
$({}_2X_S)\setminus\{0\}$ \unsure{(}car $\pm1$ opère trivial\textsuperscript{t} - donc
${}_2X_S\setminus\{0\}$ est un torseur sous
$(\mathbb Z/6\mathbb Z)/\pm1\simeq\mathbb Z/3\mathbb Z$
(\deleted{\unsure{conj.}} (générateur correspondant
$\unsure{\simeq|=}\,-\bar\zeta\,\unsure{j|\zeta}=\exp\frac{2\pi}{6}$),
\deleted{\unsure{en ce cas}} \inserted{le torseur}
s'identifie : $S/\pm1$, en \unsure{\deleted{\illegible{1}}\inserted{a}}ssociant
\unsure{à|:} chaque $s\in S$ le milieu du rayon
$[0,s]$, ou plutôt son image dans \unsure{\inserted{)}} $X_S$.

L'ens des pts fixes de $\mathbb Z/3\mathbb Z\inserted{\subset\mathbb Z/6\mathbb Z}$ sur $X_S$
est de card 3, il s'obtient en prenant
les barycentres des six triangles
de la subdivision centrale (bibarycentrique)
de l'hexagone (pts cerclés) (tels que
$(1-\bar\zeta)/3$ etc dans le cas \unsure{:} épinglé), ils
se groupent en deux \unsure{paquets|régiments} de 3
mutuellemen\unsure{t} congrus sous $\Gamma$, donc \unsure{donnant|donnent}
deux \unsure{é|e}léments (symétriques) de $X_S$, d'ordre 3,

\margin{10}
